\uuid{BLqM}
\titre{Espaces supplémentaires/Base/Dimension/Polynômes}
\auteur{hanani ; mzari}
\organisation{exo7}
\licence{CC BY-NC-SA 4.0}
\datecreate{2020-01-20}
\niveau{L1}
\module{Algèbre}
\chapitre{Espace vectoriel}
\sousChapitre{Somme directe}
\difficulte{3}
\theme{Espaces vectoriels supplémentaires, Espaces supplémentaires, Base, Dimension, Polynômes}
\exo7qcmid{lille-2.espaces-vectoriels.056}

\contenu{
\begin{question}
\qclassification{106.03}
\qtags{motcle=Espaces supplémentaires/Base/Dimension/Polynômes}
Dans $\Rr_3[X]$, l'espace des polynômes à coefficients réels de degré $\le 3$, on considère les deux sous-espaces vectoriels :
$$E=\{P\in \Rr_3[X]\; ; \; P(-X)=P(X)\}\; \mbox{ et }\; F=\{P\in \Rr_3[X]\; ; \; P(-X)=-P(X)\}.$$
Quelles sont les assertions vraies ?
\begin{answers}
\good{$\dim E= 2$}.
\bad{$\dim F = 3$}.
\bad{$\dim E\cap F = 1$}.
\good{$E$ et $F$ sont supplémentaires dans $\Rr_3[X]$}.
\end{answers}
\begin{explanations}  Une base de $E$ est $\{1,X^2\}$,  donc $\dim E=2$. Une base de $F$ est $\{X,X^3\}$, donc $\dim F=2$. On vérifie que $E\cap F  =\{0\}$. Donc, d'après le théorème de la dimension d'une somme, $\dim (E+F) = 4$ et comme $E+F$ est un sous-espace de $\Rr_3[X]$ et $\dim \Rr_3[X]=4$, $E+F= \Rr_3[X]$. Ainsi, $E$ et $F$ sont supplémentaires dans $\Rr_3[X]$.
\end{explanations}
\end{question}
}
